\chapter{Encoder-Only Transfer Across Grids}
\label{ch:encoder-transfer-bridge}

\begingroup
\color{red}

\newcommand{\plainpending}[1]{[Pending: #1]}
\newcommand{\plainpendingplot}[2][0.20\textheight]{%
  \setlength{\fboxrule}{0.7pt}%
  \fbox{\begin{minipage}[c][#1][c]{0.90\linewidth}%
    \centering\small #2%
  \end{minipage}}%
}

Screens A--E give graph encoders that work well on \texttt{bus14}. Their
action heads cannot be reused on WCCI because each output corresponds to one
action from the \texttt{bus14} action list. In this chapter, I keep the
encoder, freeze it, remove the old action heads, and learn new heads on WCCI.
Chapter~\ref{ch:screen-f-transfer} studies the other case where the encoder and
the candidate-action scorer are transferred together.

\section{Question and hypotheses}
\label{sec:encoder-transfer-question}

The main question is simple: does an encoder trained on \texttt{bus14} give
better WCCI features than an encoder that has not been trained on the source
grid?

The experiment checks three things:

\begin{itemize}
  \item source pretraining against random frozen features;
  \item original inputs against transformed inputs;
  \item a frozen encoder against the same architecture trained end to end on
  WCCI.
\end{itemize}

Four representations are kept from Screens D and E:

\begin{itemize}
  \item \texttt{e0n0v0}: the basic busbar graph;
  \item \texttt{e1n0v0}: the busbar graph with same-substation edges;
  \item \texttt{gbi\_lbi}: the disaggregated graph with bidirectional generator
  and load edges;
  \item \texttt{ga2b\_lb2a}: the disaggregated graph with asymmetric generator
  and load edges.
\end{itemize}

They all use the common choices from Screens A--C: GINE, two message-passing
layers, width 128, and energized/max graph readout. I keep several
representations because their results are very close on \texttt{bus14}, but
they may behave differently on WCCI.

\section{Transferred components and controls}
\label{sec:encoder-transfer-design}

The setup is the following:

\begin{itemize}
  \item \textbf{Frozen:} node and edge projections, type embeddings when they
  exist, the GINE layers, normalization layers, and the final representation
  projection.
  \item \textbf{Replaced:} the old \texttt{bus14} action heads.
  \item \textbf{Trained on WCCI:} one new action head per agent and a new
  centralized critic.
\end{itemize}

The four training routes are shown in Table~\ref{tab:encoder-transfer-routes}.

\begin{table}[H]
  \centering
  \small
  \caption{Training routes for the encoder-only transfer study.}
  \label{tab:encoder-transfer-routes}
  \setlength{\tabcolsep}{5pt}
  \color{red}
  \begin{tabular}{L{0.22\textwidth}L{0.25\textwidth}L{0.25\textwidth}L{0.18\textwidth}}
    \toprule
    Route & Encoder initialization & Trainable actor parameters on WCCI & Purpose \\
    \midrule
    Target scratch & Random & Encoder and new heads & Target-learning reference \\
    Random frozen & Random & New heads only & Frozen-feature control \\
    Original-input frozen transfer & Existing original-input \texttt{bus14} checkpoint & New heads only & Original-input encoder reuse \\
    Transformed-input frozen transfer & Retrained transformed-input \texttt{bus14} checkpoint & New heads only & Transformed-input encoder reuse \\
    \bottomrule
  \end{tabular}
\end{table}

The random-frozen route is important. A new WCCI head may learn something even
if its encoder is random. A source encoder only shows useful transfer if it
does better than this control. The scratch route shows what the architecture
can learn when all its weights are allowed to change on WCCI.

\begin{table}[H]
  \centering
  \small
  \caption{Protocol fields that must be fixed before the matched WCCI
  experiment. Every route must use the same target values.}
  \label{tab:encoder-transfer-protocol}
  \setlength{\tabcolsep}{6pt}
  \color{red}
  \begin{tabular}{L{0.31\textwidth}L{0.57\textwidth}}
    \toprule
    Field & Fixed value \\
    \midrule
    WCCI reduced action space & \plainpending{action cap and dataset path} \\
    WCCI training budget & \plainpending{environment steps} \\
    Source and target seeds & \plainpending{seed set} \\
    Head optimizer and learning rate & \plainpending{matched head-training recipe} \\
    Checkpoint selection & \plainpending{selection score and rollout frequency} \\
    Final evaluation & 50 fixed test chronics; 24 difficult chronics \\
    Maintenance and heuristic & Maintenance disabled; \plainpending{gate condition} \\
    \bottomrule
  \end{tabular}
\end{table}

For the original-input route, I use the existing checkpoints from Screens
A--E. For the transformed-input route, the same architectures are retrained on
\texttt{bus14} with physical power scaling and line-angle differences. The
same transformation is then used on WCCI. It cannot be added only during
transfer because the frozen encoder must receive the same kind of input that
it saw during source training.

\section{Results}
\label{sec:encoder-transfer-results}

\subsection{Learning dynamics}

The learning curves will compare:

\begin{itemize}
  \item target scratch;
  \item random frozen;
  \item original-input frozen transfer;
  \item transformed-input frozen transfer.
\end{itemize}

All routes need the same WCCI budget and the same evaluation steps. The main
thing to check is whether a source encoder learns faster or reaches a better
score than the random-frozen control.

\begin{figure}[H]
  \centering
  \plainpendingplot[0.24\textheight]{Replace with matched WCCI learning curves.
  Use the four routes above and the same axes for every representation.}
  \caption{Learning curves for the encoder-transfer routes.}
  \label{fig:encoder-transfer-learning-curves}
\end{figure}

\subsection{Full-test comparison}

The final comparison reports:

\begin{itemize}
  \item mean survival on all 50 chronics;
  \item mean survival on the 24 difficult chronics;
  \item difficult chronics completed by the model;
  \item easy chronics that remain completed.
\end{itemize}

\begin{table}[H]
  \centering
  \small
  \caption{Mean survival on the 24 difficult WCCI chronics. Values will be
  added after the matched evaluations are complete.}
  \label{tab:encoder-transfer-difficult}
  \setlength{\tabcolsep}{4.5pt}
  \color{red}
  \begin{tabular}{L{0.18\textwidth}L{0.15\textwidth}L{0.15\textwidth}L{0.15\textwidth}L{0.15\textwidth}}
    \toprule
    Representation & Target scratch & Random frozen & Original-input frozen & Transformed-input frozen \\
    \midrule
    \texttt{e0n0v0}       & TBD & TBD & TBD & TBD \\
    \texttt{e1n0v0}       & TBD & TBD & TBD & TBD \\
    \texttt{gbi\_lbi}     & TBD & TBD & TBD & TBD \\
    \texttt{ga2b\_lb2a}   & TBD & TBD & TBD & TBD \\
    \bottomrule
  \end{tabular}
\end{table}

\begin{figure}[H]
  \centering
  \plainpendingplot[0.25\textheight]{Replace with the final comparison. Show
  overall survival, difficult survival, rescues, and easy chronics kept.}
  \caption{Final WCCI comparison of the four representations and four training
  routes. Do nothing and scratch remain visible as references.}
  \label{fig:encoder-transfer-endpoints}
\end{figure}

\noindent\textbf{Result to add:}
\begin{itemize}
  \item Does a pretrained frozen encoder beat the random-frozen control?
  \item Are transformed inputs better than original inputs?
  \item How far is the best frozen encoder from target scratch?
\end{itemize}

\begin{tcolorbox}[
  colframe=red!75!black,
  colback=red!5!white,
  title=Chapter 9 Summary: Encoder-Only Transfer]
\color{red}
\begin{itemize}
  \item \textbf{Purpose:} Check if an encoder learned on \texttt{bus14} still
  helps after the action heads are replaced on WCCI.
  \item \textbf{Compared:} Four source encoders, original and transformed
  inputs, a random-frozen control, and training from scratch.
  \item \textbf{Finding:} \plainpending{main difficult-chronic result and
  learning-speed result}
  \item \textbf{Takeaway:} \plainpending{whether the encoder transfers and
  whether transformed inputs help}
\end{itemize}
\end{tcolorbox}

\endgroup
